\newif\ifstandalonechapter
\ifdefined\chapter
  \standalonechapterfalse
\else
  \standalonechaptertrue
  \begin{filecontents*}{background_refs.bib}
@article{benlloch2026,
  author  = {Benlloch-Caballero, J. and others},
  title   = {{E2E Network Slicing for Enhanced Cybersecurity Orchestration, Automation and Response in 5G/6G: The RIGOUROUS Approach}},
  journal = {Journal of Network and Systems Management},
  year    = {2026},
  volume  = {34},
  number  = {22}
}

@article{breiman2001rf,
  author  = {Breiman, L.},
  title   = {{Random Forests}},
  journal = {Machine Learning},
  volume  = {45},
  number  = {1},
  pages   = {5--32},
  year    = {2001}
}

@inproceedings{chen2016xgboost,
  author    = {Chen, T. and Guestrin, C.},
  title     = {{XGBoost: A Scalable Tree Boosting System}},
  booktitle = {Proc. ACM SIGKDD},
  year      = {2016},
  pages     = {785--794}
}

@article{dawn2025,
  author  = {Lalli, P. and Olmsted, A. and Tripathi, S.},
  title   = {{DAWN Framework: SDN, NFV, and Machine Learning for DDoS Resistance}},
  booktitle = {Lecture Notes in Computer Science},
  year    = {2025},
  publisher = {Springer}
}

@inproceedings{deoliveira2023,
  author    = {de Oliveira, R. L. S. and others},
  title     = {{Optimal VNF Placement for DDoS Mitigation in Cloud Data Centers}},
  booktitle = {Proc. IEEE CLOUD},
  year      = {2023},
  pages     = {120--130}
}

@techreport{etsi_nfv_arch,
  author      = {{ETSI NFV ISG}},
  title       = {{Network Functions Virtualisation: Architectural Framework}},
  institution = {ETSI},
  type        = {GS NFV 002},
  year        = {2014}
}

@article{jcc2026proactive,
  author  = {{Various Authors}},
  title   = {{Online Machine Learning for Proactive Threat Intelligence and Timely DDoS Attack Prediction}},
  journal = {Journal of Cloud Computing},
  year    = {2026}
}

@article{jnsm2025qosentry,
  author  = {{Various Authors}},
  title   = {{QoSentry: A Reinforcement Learning Framework for QoS-Preserving DDoS Mitigation in SDN}},
  journal = {Journal of Network and Systems Management},
  year    = {2025}
}

@article{ke2017lightgbm,
  author  = {Ke, G. and others},
  title   = {{LightGBM: A Highly Efficient Gradient Boosting Decision Tree}},
  journal = {Advances in Neural Information Processing Systems},
  volume  = {30},
  year    = {2017}
}

@inproceedings{khozam2024,
  author    = {Khozam, A. and Blanc, G. and Tixeuil, S. and Totel, E.},
  title     = {{DAMS: Double Deep Q-Network Approach for DDoS Attack Mitigation in SDN}},
  booktitle = {Proc. ICFNDS},
  year      = {2024}
}

@article{lundberg2017unified,
  author  = {Lundberg, S. M. and Lee, S.-I.},
  title   = {{A Unified Approach to Interpreting Model Predictions}},
  journal = {Advances in Neural Information Processing Systems},
  volume  = {30},
  year    = {2017}
}

@article{mdpi2025transformer,
  author  = {{Various Authors}},
  title   = {{A Transformer-Based Framework for DDoS Attack Detection via Temporal Dependency and Behavioral Pattern Modeling}},
  journal = {Algorithms},
  year    = {2025},
  volume  = {18},
  number  = {10}
}

@article{mirkovic2004taxonomy,
  author  = {Mirkovic, J. and Reiher, P.},
  title   = {{A Taxonomy of DDoS Attack and DDoS Defense Mechanisms}},
  journal = {ACM SIGCOMM Computer Communication Review},
  volume  = {34},
  number  = {2},
  pages   = {39--53},
  year    = {2004}
}

@misc{onap_arch,
  author       = {{ONAP Community}},
  title        = {{ONAP Architecture Overview}},
  year         = {2024},
  howpublished = {\url{https://docs.onap.org/en/latest/guides/onap-user/architecture.html}}
}

@article{scifrep2025cnngru,
  author  = {{Various Authors}},
  title   = {{DDoS Classification of Network Traffic in Software Defined Networking SDN Using a Hybrid Convolutional and Gated Recurrent Neural Network}},
  journal = {Scientific Reports},
  year    = {2025}
}

@inproceedings{sdnshield2021,
  author    = {{Various Authors}},
  title     = {{SDNShield: NFV-Based Defense Framework Against DDoS Attacks on SDN Control Plane}},
  booktitle = {Proc. IEEE TNSM},
  year      = {2021}
}

@inproceedings{vaswani2017attention,
  author    = {Vaswani, A. and others},
  title     = {{Attention Is All You Need}},
  booktitle = {Advances in Neural Information Processing Systems},
  volume    = {30},
  year      = {2017}
}
  \end{filecontents*}
  % !TeX program = xelatex
  \documentclass[12pt, a4paper, reqno, oneside]{book}
  \usepackage[left=2.5cm, right=2cm, top=2cm, bottom=2cm]{geometry}
  \usepackage{iftex}
  \ifPDFTeX
    \usepackage[utf8]{inputenc}
    \usepackage[T5]{fontenc}
    \usepackage[vietnamese,english]{babel}
    \usepackage{vntex}
  \else
    \usepackage{fontspec}
    \usepackage{polyglossia}
    \setmainlanguage{vietnamese}
    \setotherlanguage{english}
    \setmainfont{Times New Roman}
  \fi
  \usepackage{amsmath,amsxtra,amssymb,latexsym, amscd,amsthm}
  \usepackage{microtype}
  \usepackage{indentfirst}
  \usepackage{color}
  \usepackage{empheq}
  \usepackage{amsmath}
  \usepackage[mathscr]{eucal}
  \usepackage{amsfonts}
  \usepackage{graphicx}
  \usepackage{float}
  \usepackage{listings}
  \usepackage[dvipsnames]{xcolor}
  \usepackage{verbatim}
  \usepackage{tikz}
  \usetikzlibrary{calc}
  \usepackage{tkz-graph}
  \usetikzlibrary{positioning}
  \usepackage{float}
  \usepackage{array}
  \usepackage{multirow}
  \usepackage{algorithm}
  \usepackage{algpseudocode}
  \usepackage{subfigure}
  \usepackage{pgfplots}
  \usepackage{titlesec}
  \usepackage{fancyhdr}
  \usepackage{rotating}
  \usepackage{longtable}
  \usepackage{supertabular}
  \fancyhf{}
  \fancyhead[LO,RE]{\small{\leftmark}}
  \fancyfoot[C]{\thepage}
  \usepackage{ulem}
  \usepackage{pdfpages}
  \usepackage{niceframe}
  \usepackage[square,numbers]{natbib}
  \usepackage{subcaption,booktabs}
  \usepackage{chngcntr}
  \usepackage{tabularray}
  \usepackage[labelfont=bf, textfont=it, justification=centering]{caption}
  \captionsetup{labelfont={color=blue, bf}, textfont={color=black, it}}
  \renewcommand{\labelitemi}{+}
  \hyphenation{op-tical net-works semi-conduc-tor}
  \usepackage[hyphens,spaces,obeyspaces]{url}
  \usepackage[colorlinks,bookmarksopen,bookmarksnumbered,citecolor=blue, colorlinks=true,linkcolor=blue,urlcolor=blue]{hyperref}
  \usepackage{tabularx}
  \pagestyle{fancy}
  \begin{document}
\fi

\section{Phân loại và Đặc trưng luồng mạng của tấn công DDoS}
\label{sec:ddos_taxonomy}

Để xây dựng một hệ thống phát hiện và giảm thiểu hiệu quả, việc trích xuất các đặc trưng tĩnh và động của lưu lượng mạng (flow-based features) là cực kỳ quan trọng~\cite{mirkovic2004taxonomy}. Dựa trên các hành vi mạng ở tầng giao thức (Layer 3/4) và tầng ứng dụng (Layer 7), các cuộc tấn công DDoS hiện đại được phân loại thành 6 nhóm chính, mỗi nhóm thể hiện các dấu vết đặc trưng riêng biệt trên dữ liệu NetFlow/sFlow như được tóm tắt trong Bảng~\ref{tab:ddos_classification}.

\begin{table}[H]
\centering
\caption{Phân loại các dạng tấn công DDoS dựa trên đặc trưng luồng mạng}
\label{tab:ddos_classification}
\begin{tabularx}{\textwidth}{|c|l|X|}
\hline
\textbf{ID} & \textbf{Loại tấn công} & \textbf{Đặc trưng lưu lượng (Flow-based Features)} \\ \hline
1 & \textbf{UDP Flood} & Hoạt động dựa trên nguyên lý vắt kiệt băng thông. Dấu hiệu: tốc độ gói tin (\texttt{pkt\_rate}) rất cao, phân phối giao thức thiên về UDP (\texttt{proto\_dist\_udp} cao), và kích thước gói tin trung bình (\texttt{avg\_pkt\_size}) thấp. \\ \hline
2 & \textbf{SYN Flood} & Lợi dụng cơ chế bắt tay 3 bước của TCP. Dấu hiệu: tỷ lệ cờ SYN (\texttt{syn\_ratio}) đột biến, tốc độ sinh luồng mới (\texttt{new\_flows\_rate}) rất cao nhưng thời gian tồn tại trung bình của mỗi luồng lại rất thấp (\texttt{flow\_duration\_mean}). \\ \hline
3 & \textbf{HTTP Flood} & Tấn công bạo lực ở Layer 7. Dấu hiệu: \texttt{new\_flows\_rate} cực kỳ cao do botnet liên tục gửi yêu cầu GET/POST, chạy trên nền tảng TCP (\texttt{proto\_dist\_tcp} $\approx 1$), và \texttt{avg\_pkt\_size} nằm trong mức trung bình (500–1400B). \\ \hline
4 & \textbf{ICMP Flood} & Dấu hiệu cơ bản là \texttt{proto\_dist\_icmp} $\approx 1$. Đặc biệt, độ hỗn loạn của IP nguồn (\texttt{src\_ip\_entropy}) sẽ ở mức rất cao nếu kẻ tấn công sử dụng kỹ thuật giả mạo IP (IP Spoofing) để né tránh các luật chặn thông thường. \\ \hline
5 & \textbf{Amplification} & Kỹ thuật khuếch đại (DNS, NTP). Dấu hiệu: tốc độ dữ liệu (\texttt{byte\_rate}) cao bất thường, kích thước gói tin cực lớn (\texttt{large avg\_pkt\_size}), trong khi tốc độ tạo gói tin (\texttt{pkt\_rate}) lại tương đối thấp so với lưu lượng. \\ \hline
6 & \textbf{Slow-rate} & Tấn công tàng hình Layer 7 (như Slowloris). Dấu hiệu đối lập hoàn toàn với các dạng bạo lực: thời gian luồng (\texttt{flow\_duration\_mean}) kéo dài tối đa, trong khi \texttt{byte\_rate} và \texttt{pkt\_rate} cực thấp, \texttt{proto\_dist\_tcp} $\approx 1$. \\ \hline
\end{tabularx}
\end{table}

Sự đa dạng về đặc trưng này—từ các cuộc tấn công bạo lực tiêu tốn băng thông (Volumetric) cho đến các cuộc tấn công tàng hình tốc độ thấp (Slow-rate)—đặt ra thách thức lớn cho các hệ thống phòng thủ dựa trên ngưỡng (threshold-based) truyền thống, đòi hỏi sự tham gia của các thuật toán Học máy và khả năng điều phối mạng tự động.

\section{Research Gap}
\label{sec:research_gap}

Tổng quan tài liệu cho thấy khoảng trống nghiên cứu nằm ở giao điểm của ba yêu cầu: phát hiện chủ động, điều phối vòng kín theo chuẩn công nghiệp, và khả năng giải thích quyết định tự động.

\begin{enumerate}
  \item \textbf{Thiếu cơ chế chủ động:} Phần lớn hệ thống hiện nay vẫn phản ứng sau khi tấn công đã xảy ra, làm tăng độ trễ phản ứng đầu-cuối~\cite{mirkovic2004taxonomy}.
  \item \textbf{Thiếu tích hợp MANO/ONAP thực tế:} Nhiều mô hình có kết quả tốt ở lớp phát hiện nhưng chưa nối chặt với vòng điều phối thực thi trên hạ tầng chuẩn~\cite{etsi_nfv_arch,onap_arch}.
  \item \textbf{Thiếu đo lường end-to-end và explainability vận hành:} Ít công trình báo cáo đầy đủ pipeline phát hiện-dieu phoi-thuc thi kèm khả năng audit quyết định~\cite{benlloch2026}.
\end{enumerate}

\section{Problem Statement}
\label{sec:problem_statement}

Hệ thống phòng thủ DDoS hiện nay trong môi trường cloud-native và 5G/6G vẫn đối mặt bốn vấn đề cốt lõi.

Thứ nhất, cơ chế phát hiện còn mang tính phản ứng: phần lớn giải pháp chỉ kích hoạt sau khi lưu lượng tấn công đã bùng phát, làm tăng độ trễ phản ứng đầu cuối và giảm khả năng bảo toàn dịch vụ~\cite{jcc2026proactive}.

Thứ hai, nhiều giải pháp học máy chưa được tích hợp chặt với lớp điều phối mạng chuẩn công nghiệp. Điều này tạo khoảng cách giữa kết quả mô hình và hành động thực thi trong hạ tầng MANO/ONAP~\cite{etsi_nfv_arch,onap_arch}.

Thứ ba, bài toán phát hiện trong thực tế phức tạp hơn do tấn công đa vectơ, tấn công tốc độ thấp và đặc trưng lưu lượng biến thiên theo thời gian (đã được trình bày tại Bảng~\ref{tab:ddos_classification}). Các mô hình đơn lẻ thường khó đồng thời đạt độ chính xác cao, độ trễ thấp và tính ổn định khi dữ liệu mất cân bằng~\cite{mdpi2025transformer}.

Thứ tư, phần lớn hệ thống tự động thiếu khả năng giải thích quyết định ở mức vận hành, gây khó khăn cho kiểm chứng và audit trước khi thực thi chính sách chặn hoặc chuyển hướng lưu lượng~\cite{lundberg2017unified}.

Từ các vấn đề trên, luận văn tập trung vào mục tiêu xây dựng một pipeline phát hiện và giảm thiểu DDoS có thể: (i) cảnh báo sớm, (ii) tích hợp trực tiếp với vòng kín ONAP, (iii) đo lường end-to-end, và (iv) cung cấp thông tin giải thích phục vụ vận hành.

\section{Related Works}
\label{sec:related_works}

Các công trình liên quan có thể chia thành bốn hướng chính.

\textbf{(1) SDN/NFV-based defense frameworks.} DAWN và SDNShield cho thấy lợi ích của ảo hóa chức năng mạng trong giảm thiểu DDoS, nhưng phạm vi triển khai còn cục bộ và chưa gắn đầy đủ với orchestration chuẩn~\cite{dawn2025,sdnshield2021}.

\textbf{(2) Reinforcement Learning cho quyết định giảm thiểu.} QoSentry và DAMS chứng minh RL có thể tối ưu trade-off giữa QoS và hiệu quả phòng thủ; tuy nhiên chi phí huấn luyện, hội tụ chính sách và khả năng triển khai thực tế vẫn là thách thức~\cite{jnsm2025qosentry,khozam2024}.

\textbf{(3) ML/DL cho phát hiện và dự báo DDoS.} Nhóm mô hình truyền thống (Random Forest, XGBoost, LightGBM) mạnh về tốc độ suy luận; nhóm DL/Transformer mạnh về biểu diễn phụ thuộc thời gian nhưng thường đánh đổi bằng độ trễ cao hơn~\cite{mdpi2025transformer}.
\begin{table}[h]
\centering
\caption{Thống kê các kết quả SOTA trên bộ dữ liệu CICDDoS2019}
\label{tab:cicddos2019_sources}
\begin{tabular}{|l|c|p{8cm}|}
\hline
\textbf{Phương pháp} & \textbf{Độ chính xác} & \textbf{Tên bài báo / Nguồn gốc thông tin} \\ \hline
Random Forest & 99.9974\% & Detecting Distributed Denial of Service Attacks (Hudaib et al.) \\ \hline
ELM + Black-hole & 99.80\% & Optimized extreme learning machine for detecting DDoS attacks in cloud computing (Kamber et al.) \\ \hline
CNN-LSTM Lai & 99.63\% & Empirical evaluation of ensemble learning and hybrid CNN-LSTM for IoT threat detection (Ahmed et al.) \\ \hline
TD3 (DRL) & 99.12\% & Cloud-based DDoS detection using hybrid feature selection with deep reinforcement learning (Satpathy et al.) \\ \hline
Standalone LSTM & 99.40\% & Evaluating the Performance of LSTM and GRU in Detection of DDoS Attacks (Subramanian et al.) \\ \hline

\end{tabular}
\end{table}

\textbf{(4) ONAP-integrated closed-loop security.} RIGOUROUS là hướng tiếp cận gần nhất với triển khai công nghiệp khi tích hợp ONAP và network slicing đầu-cuối. Dù vậy, điểm hạn chế chính vẫn là thiên về phản ứng sau khi tấn công xuất hiện và chưa tích hợp sâu thành phần giải thích quyết định~\cite{benlloch2026,deoliveira2023}.

Tổng quan cho thấy khoảng trống nghiên cứu tập trung ở giao điểm giữa \textit{proactive detection}, \textit{closed-loop orchestration} và \textit{explainability}.

\begin{table}[H]
\centering
\caption{Tổng hợp so sánh công trình liên quan và khoảng trống}
\label{tab:related_gap}
\begin{tabularx}{\textwidth}{lXccccc}
\toprule
\textbf{Công trình} & \textbf{Phương pháp} & \textbf{Chủ động} & \textbf{MANO thực} & \textbf{E2E đo} & \textbf{XAI} & \textbf{Tài liệu} \\
\midrule
RIGOUROUS (2026)   & RL + ONAP      & \texttimes & \checkmark & \texttimes & \texttimes & \cite{benlloch2026} \\
DAWN (2025)        & SDN+NFV+ML     & \texttimes & \texttimes & \texttimes & \texttimes & \cite{dawn2025} \\
QoSentry (2025)    & RL + SDN       & \texttimes & \texttimes & \texttimes & \texttimes & \cite{jnsm2025qosentry} \\
DAMS (2024)        & DQN + SDN      & \texttimes & \texttimes & \texttimes & \texttimes & \cite{khozam2024} \\
de Oliveira (2023) & VNF Placement  & \texttimes & \texttimes & \texttimes & \texttimes & \cite{deoliveira2023} \\
\midrule

\bottomrule
\end{tabularx}
\end{table}

\ifstandalonechapter
\bibliographystyle{IEEEtran}
\bibliography{background_refs}
\end{document}
\fi